\UTIResumenPagina
  {english}
  {Abstract}
  {Ecuador's Organic Law on Personal Data Protection requires every data controller to use a risk analysis and management methodology, and failing to do so is a serious infringement. In Ecuador, 16,161 educational institutions and 280 higher education institutions are subject to this obligation, while the data protection officers registered in April 2025 amounted to at most 7.0\,\% of them, and there is no official measurement of universities' compliance. The general objective was to design a methodology for diagnosing and managing risks in the processing of personal data of students and faculty in higher education institutions, integrating the risk management guide of the Superintendency of Personal Data Protection with ISO/IEC 27005:2022. Design science research with a mixed approach was followed, applied to the admission and enrollment process of Universidad Tecnológica Indoamérica, Quito campus. The methodology comprises a preparation phase, a maturity diagnosis phase and four management phases (identification, assessment, mitigation and monitoring), mapped to ISO/IEC 27035-1:2023 and NIST SP 800-61r3; it estimates annual likelihood and impact on data subjects through a Delphi panel calibrated with the Brier score, and assigns risk levels with a decision table that prevents severe harm from being hidden behind low likelihood. As verified results, ten requirements for performing the diagnosis, thirteen candidate scenarios traceable to the law, and a matrix template whose 1,057 formulas evaluate without errors and which passes six test cases were obtained. The results of the institutional application and of expert validation depend on institutional authorization and are identified as pending. It is concluded that the methodology is fully specified and that its validity will be judged against twelve indicators with targets defined in advance.}
  {Keywords}
  {data protection impact assessment, higher education institutions, personal data protection, risk management}
